% =====================================================================
% RESULTADOS QUE APORTA EL EQUIPO
% Este es el único archivo que hay que editar después de las pruebas con
% usuario y de repetir las herramientas. Cada valor aparece en todas las
% secciones que lo usan. Reemplace el texto entre llaves del \newcommand.
% =====================================================================

% ---------------------------------------------------------------------
% 1. Prueba de aceptación (Segunda parte, sección 6; Apéndice G)
% ---------------------------------------------------------------------
\newcommand{\resUsuarioFinal}{Gonzalo Nicolas Morales Ciriza, usuario externo al equipo}
\newcommand{\resFechaSesion}{25/11/2026}
\newcommand{\resCommitDesplegado}{la rama \codigo{main} después de integrar \codigo{fix/defectos}}
\newcommand{\resCasosEjecutados}{17 de 17}
\newcommand{\resCasosAceptados}{17 aceptados y 0 rechazados}
\newcommand{\resRFAceptados}{20/20 = 100\,\%}
\newcommand{\resObservaciones}{(1) «La primera vez que abrí la página en el celular tardó en cargar»; (2) «Después de un rato sin usarla, el resumen tardó unos segundos en mostrar los datos»; (3) «Sería útil que la última categoría que usé apareciera primero al registrar un gasto»}
\newcommand{\resAcciones}{(1) mejora futura: bajar el LCP en perfil móvil (Tercera parte, eficiencia de desempeño); (2) mitigada: es D-09, con el aviso «Despertando el servidor» y el ping programado a \codigo{/salud}; (3) mejora futura: se agrega al backlog como historia de usuario}
\newcommand{\resConcepto}{Aceptado con observaciones}
\newcommand{\resActa}{\IfFileExists{images/firma_usuario.png}{diligenciada y firmada digitalmente}{\pendiente{firma digital del usuario final en images/firma\_usuario.png}}}

% ---------------------------------------------------------------------
% 2. Tarea cronometrada PU-06 (Segunda parte, 5.3; Tercera parte, usabilidad)
% ---------------------------------------------------------------------
\newcommand{\resPUParticipantes}{Leonel Garrido Naeder, Edimes Amaya Ballestas, Sabrina Ferreira Ferreira y Gonzalo Nicolas Morales Ciriza, usuarios externos al equipo, e Isabella Amaya, integrante del equipo}
\newcommand{\resPUMediana}{36}
\newcommand{\resPUExito}{5/5 = 1,00}
\newcommand{\resPUObservaciones}{tiempos de 24, 33, 36, 41 y 52 s; dos participantes buscaron primero la categoría «Alimentación» antes de escribir el monto y ninguno pidió ayuda}

% ---------------------------------------------------------------------
% 3. Herramientas que se pueden repetir en vivo (por defecto, el texto
%    dice que no se repitieron; si las corren, reemplazar por el resultado)
% ---------------------------------------------------------------------
\newcommand{\resZapApi}{26/11/2026: 0 altas, 1 media (respuesta 500 con contraseñas de más de 72 bytes, límite de bcrypt; hallazgo abierto), 4 bajas, 116 reglas superadas}
\newcommand{\resZapFrontend}{26/11/2026: 0 altas; las 3 medias iniciales quedaron corregidas y aparecen 6 medias que piden endurecer la CSP (\codigo{unsafe-inline}, comodín), SRI y CORS; 4 bajas; 57 reglas superadas}
\newcommand{\resZapResumen}{0 / 6 (26/11/2026)}
\newcommand{\resSonar}{26/11/2026, Community 26.11: Quality Gate aprobado; seguridad A (0), mantenibilidad A (63), fiabilidad D (110: 103 del frontend, sobre todo reglas de accesibilidad y estilo); cobertura 95,4\,\% en \codigo{backend} y 0\,\% en \codigo{frontend/src} (37,8\,\% global); duplicación 0,5\,\%}
\newcommand{\resSonarLoc}{11\,000 aprox. (26/11/2026)}
\newcommand{\resJmeterRender}{no se ejecutó contra Render porque la API no respondía el 21/11/2026}
\newcommand{\resLambdaEscritorio}{no verificados (se agotaron las sesiones gratuitas de LambdaTest)}
\newcommand{\resLambdaMovil}{no verificados en dispositivos reales}
\newcommand{\resNavegadores}{2/4 = 0,50 (Edge y Safari sin verificar)}
\newcommand{\resDisponibilidadApi}{la API dejó de responder porque Supabase pausó la base por inactividad; el 25/11/2026, al reactivarla, volvió a responder. Desde entonces el ping llama a \codigo{/salud}, que también consulta la base}
\newcommand{\resPIcincoPostgres}{no ejecutado sobre PostgreSQL; verificado en SQLite con las claves foráneas activadas}
